\scanpage{55}
\mspage{35}

c) Finally, let $\rho_F:F\longrightarrow P_XF$ be the presheaf morphism
such that $(\rho_F)_U:F(U)\longrightarrow P_XF(U)$ is the unique morphism with
$\tilde{p}_{xU}\circ(\rho_F)_U=r_{xU}$ for all $x\in U$. The diagram
\diagram{\begin{tikzcd}[ampersand replacement=\&,column sep=large,row sep=large]
F(U) \arrow[r,"r_{VU}"] \arrow[d,"(\rho_F)_U"'] \&
F(V) \arrow[r,"r_{xU}"] \arrow[d,"(\rho_F)_V"'] \&
F_x \arrow[d,"i_{F_x}"']\\
P_XF(U) \arrow[r,"p_{VU}"'] \&
P_XF(V) \arrow[r,"\tilde{p}_{xV}"'] \& F_x
\end{tikzcd}}
is then commutative. Taking direct limits of the second square
yields a commutative diagram
\diagram{\begin{tikzcd}[ampersand replacement=\&,column sep=large,row sep=large]
F(U) \arrow[r,"r_{xU}"] \arrow[d,"(\rho_F)_U"'] \&
F_x \arrow[d,"(\rho_F)_x"'] \arrow[dr,"i_{F_x}"] \& \\
P_XF(U) \arrow[r,"p_{xU}"'] \&
(P_XF)_x \arrow[r,"\tilde{p}_x"'] \& F_x
\end{tikzcd}}

The $\rho_F$ are the components of a natural transformation
$\rho:I_{\underline{P}(X,\underline{A})}\longrightarrow P_X$; i.e., for any $\varphi:F\longrightarrow F'$,
$(\rho_{F'})\circ\varphi=P_X\varphi\circ\rho_F$. It is sufficient to check this on an
open set $U$, and, in fact, it is sufficient to show that
\[
\tilde{p}'_{xU}\circ(\rho_{F'})_U\circ\varphi_U
=\tilde{p}'_{xU}\circ(P_X\varphi)_U\circ(\rho_F)_U
\]
But
\[
\begin{aligned}
\tilde{p}'_{xU}\circ(\rho_{F'})_U\circ\varphi_U
&=r'_{xU}\circ\varphi_U=\varphi_x\circ r_{xU}=\\
&\quad\varphi_x\circ\tilde{p}_{xU}\circ(\rho_F)_U
=\tilde{p}'_{xU}\circ(P_X\varphi)_U\circ(\rho_F)_U
\end{aligned}
\]
\mspage{54}
